Elke dag bij haar volk
hing een donkere wolk

boven Maria’s hoofd.
Zo had ze nooit geloofd,

dat ze kalm kon spelen
en haar ziel kon delen.

Ze kon worden gestoord,
beschuldigd en verhoord.

Zo ging ze geloven
dat een rechter boven

haar hoofd zijn schaduw wierp.
En haar bestaan verwierp.

Zo ontstond haar waakstand:
“Hou alles in de hand!”

Ze leek in dit kader
vrij sterk op haar vader.

Als een heilig boontje
kwam Maria’s loontje

pas uit haar kloostergang.
Dan was ze niet meer bang,

meende ze incorrect.
Haar angst werd slechts bedekt.

Koud smeltwater stroomde,
terwijl ze dagdroomde,

door beekjes naar het dal.
De landbouwstroken, smal

middels stapels stenen
afgebakend, verschenen

langs een heuvelhelling.
Kruisen ter bescherming

tegen boze geesten,
alsook wilde beesten.

Maria liep galant
door het hoge bergland

op haar kloosterroute.
Ze rook de zwoelzoete

geur van oude druiven.
Boeren wilden wuiven,

maar bedachten zich snel.
Het leidde tot een rel

over eerbied en respect,
eerder op dit traject.

Bij Mariabeeldjes,
de kleine juweeltjes

onder een afdakje,
stond een houten bakje,

voor de offers bedoeld.
Nadat ze had gevoeld

of er wel wat in zat -
men wel geschonken had -

haalde ze alles leeg.
Zodat de kerk het kreeg.

Een houten brug verbond,
vlak boven een afgrond,

het dorp met de akkers.
Enkele houthakkers

passeerden Maria
en wandelden via

de brug naar de bossen,
samen met hun ossen.

Elk liturgisch drama
was groots voor Maria.

Juist de slechte tijden
kon haar vaak verblijden.

“Val me niet steeds lastig”,
zei ze dan standvastig,

als iemand met sores
aankwam aan haar adres

op weg naar het kerkspel.
“U kent toch zeker wel

Thessalonicenzen?
Dat zijn ook mijn wensen!

U kunt mij toch geven
een rustiger leven

door uw eigen zaken
uw eigen te maken?

Stop dus uw tragedie!”,
hoorde daarna Emmy.

Ze riep hen nog eens na
op weg naar het drama.

Fel klonk het klokkenspel.
In de kleine kapel

reciteerde ze Psalm
91 kalm.

Haar woorden klonken schuw.
“… vernacht in de schaduw

van de Almachtige.”
De zenuwachtige

zag plots een hersenschim.
Het silhouet scheen slim,

maar met een slinkse blik.
“Oh mijn God, dat ben ik!

Mijn zondige natuur
schijnt zwart op deze muur.

U heeft mij doorgelicht!
Verbannen uit het licht!”

Het schilderde haar wee,
als bij Dorian Gray,

alleen in het schijnsel.
Het plotse verschijnsel

bewoog toen houterig,
gedroeg zich stekelig.

De circusartiesten,
die de schimmen spietsten,

speelden met nep gegil
blijkbaar khayal al-zill

of schaduwtheater.
Ze vatte haar flater.

De Memento Mori
was de meditatie

die Maria veel deed,
omdat ze aldoor leed.

Ze wenste te sterven,
door de vele scherven,

die trots haar vroom leven,
kervend achter bleven.

Ze stelde de dood voor,
maar die gaf geen gehoor.

De zonde van wanhoop,
zelfmoord als de afloop,

beging een doodzonde.
Slechts een gebedsstonde

bood haar wat leniging.
Hoewel haar pijniging

voortkwam uit christenstress.
Ze was een zondares

en schaamde zich daarvoor.
Maria was bang voor

het steeds tekort schieten.
Ze zag de limieten

van haar menselijkheid.
Dat, intenties ten spijt,

ze onbedoeld kwaad deed.
De mafketel verweet

zichzelf dat ze zwart zag.
Haar zelfaanvalsverdrag

was het vrome verbond.
Vrees voor dat als ze opstond,

de duivel uit haar doos
opveerde bang en boos.

Zoals Jezus ooit zei:
“Oordeel niet, opdat gij

zelf niet geoordeeld wordt.”
Doet men zichzelf tekort

dan, volgens die maatstaf,
keurt men ook iemand af

met dezelfde kenmerk.
Zou ze het schaduwwerk,

ontwikkeld door Carl Jung,
in plaats van teniet doen,

hebben geemployeerd
dan had ze dit geleerd.

In “Angst vreet de ziel op”
merkte Fassbinder op

dat vrees de liefde remt
en relaties beklemd

mits het ongezien blijft
en ons gedrag aandrijft.

Het was bij haar het wrangst
dat overlevingsangst,

terwijl ze dood wilde
haar compassie kilde.

“De grimmige wensen
van andere mensen,

gaan ten koste van mij”,
was wat haar ego zei.

Haar behoeftes moesten
anderen verwoesten.

Maria zag met spijt
haar ongevoeligheid,

maar geenszins het ontzag
dat er slinks onder lag.

Want dan zou ze weldoen
in plaats van alleen doen.

Dan liet ze het leven
om echt te gaan leven.

Dan werd Maria lief
in plaats van reactief.

Dan zag men kwetsbaarheid
in plaats van een verwijt.

Dan ervoer men een band
en niet haar harde hand.

Ze zag in Maria,
een van haar stadia,

die ze eerder doorliep.
Toen ze nog dieper sliep.

Haar vroegere versie
voldeed niet aan Emmy,

die wel haar schaduwvrees
als omkleedsel aanwees.

Zij was wakker geschud.
Elke spook heeft zijn nut.

Kon ze haar vergeven
en liet ze dit leven

weer aan zich voorbij gaan?
Zal haar schim verder gaan?

Was dit slechts het intro?
In Dante’s ‘Inferno’

ontmoet de held Dante
schaduwen of 'ombre',

zielen van de doden,
maar ook van de noden

van Dante’s eigen geest.
Had ze haar hel gevreesd?

Maakte ze een koersfout
in haar donkere woud?

Verzonk door haar dwaling
haar ziel in verwarring?

Kwam haar hel aangestormd?
Was Emmy zelf vervormd,

getekend aan zwakte?
De leegte haar pakte,

zoals Maria deed
in haar vlucht voor leed?
